\chapter*{\hfill Résumé \hfill}

Ce mémoire présente la conception et le développement de \textbf{HayaBook}, une plateforme mobile de réservation de services à la demande destinée au marché algérien. La plateforme met en relation des clients avec des prestataires de services (coiffeurs, médecins, tuteurs, etc.) via une application mobile Flutter, un tableau de bord d'administration React, et un serveur backend Node.js/Express avec une base de données PostgreSQL gérée par l'ORM Prisma. Le système intègre une authentification sécurisée par JWT, une communication en temps réel via Socket.io pour les notifications et la messagerie instantanée, ainsi qu'un moteur de réservation avec vérification de disponibilité et transitions automatiques de statuts. Les résultats obtenus démontrent la viabilité d'une solution numérique complète, scalable et sécurisée pour la gestion de services à la demande en Algérie.

\textbf{Mots-clés :} Réservation en ligne, Application mobile, Flutter, Node.js, Prisma, PostgreSQL, Socket.io, API REST, JWT, Temps réel.

